% =====================================================================
% METHODS
% =====================================================================

\section{Decomposition of the MLE test loss}
\label{sec:methods}

This section develops the analytic object the rest of the paper studies. We fix
notation and write the test loss of the maximum-likelihood plug-in as a single KL
quantity (Section~\ref{sec:setup}), introduce reference quantities that will
appear in the decomposition (Section~\ref{sec:refobj}), and state the inequality
that decomposes test loss into approximation, batch-regret, prior-cost, and MLE-vs-Bayes
contributions (Section~\ref{sec:decomp}). Section~\ref{sec:gibbs} describes how
we approximate a Bayes predictor in practice using SGLD.
Appendix~\ref{app:interpretation} summarizes the notation and the quantities
measured in the experiments.

\subsection{Setup}
\label{sec:setup}

We consider a measurable space $(\mathcal{X},\mu)$ and a true data-generating
distribution $P^*\in\mathcal{P}(\mathcal{X})$ with density $p^*=dP^*/d\mu$.
We draw training data $X_1,\dots,X_n$ and a fresh test point $X_{n+1}$ i.i.d.\
from $P^*$. A model class is a family
$\{P_{X\mid\theta}:\theta\in\Theta\}$ on $(\mathcal{X},\mu)$ with densities
$p_\theta$, which we extend to product form
$P_{X^n\mid\theta}=\prod_{i=1}^n P_{X\mid\theta}$. We make no regularity,
identifiability, smoothness, or finite-dimensionality assumption on $\Theta$, and
we allow $P^*$ to lie outside the class. We take the maximum-likelihood
estimator
\begin{equation*}
\hat\theta_n=\hat\theta_n(X^n)\in\arg\max_{\theta\in\Theta}\prod_{i=1}^n p_\theta(X_i)
\end{equation*}
to be a measurable selection, breaking ties arbitrarily. We assume the
entropies, risks, and differences below are finite. The same decomposition holds for any
measurable plug-in rule, with expectations also taken over its internal
randomness when applicable.

\paragraph{Test loss.}
We define the MLE plug-in predictor as the density
$\hat p_n\triangleq p_{\hat\theta_n(X^n)}$ and write its test log-loss as
\begin{equation}
L_n(P^*) \;\triangleq\; \mathbb{E}_{X^n,X_{n+1}\stackrel{\textit{i.i.d.}}{\sim}P^*}\!\big[\log 1/\hat p_n(X_{n+1})\big].
\label{eq:testloss}
\end{equation}
For other predictive rules $Q$, we write $L^Q_n(P)$ for the analogous loss.
Subtracting the entropy $H(P^*)\triangleq-\mathbb{E}_{P^*}[\log p^*(X)]$ and applying
the chain rule of KL divergence gives the identity
\begin{equation}
\mathcal{R}_n(P) \;\triangleq\; L_n(P)-H(P)
\;=\; \mathbb{E}_{X^n\sim P^{\otimes n}}\!\Big[D\!\left(P\,\big\|\,p_{\hat\theta_n(X^n)}\right)\Big],
\quad P\in\mathcal{P}(\mathcal{X}),
\label{eq:plugin-risk}
\end{equation}
which we call the \emph{plug-in risk} of MLE under $P$. We define
$\mathcal{R}_n(\cdot)$ for every $P$, in or out of the class; it depends on
$\Theta$ only through the MLE.

\paragraph{Supervised learning.}
Our setup generalizes immediately to supervised data, by reducing a supervised
problem to an unsupervised one with a fixed input marginal. Formally:
for $(X_i,Y_i)\stackrel{\textit{i.i.d.}}{\sim}P^*_{XY}$, we fix the
$X$-marginal $P_X=P_X^*$ and let the model
parameterize the joint
$P_{XY\mid\theta}(x,y)=P_X(x)\,P_{Y\mid X,\theta}(y\mid x)$. Substituting
$\tilde X_i\triangleq(X_i,Y_i)$ recovers the present setup; the constant
$H(P_X)$ cancels from both sides of the decomposition stated below; and the
resulting MLE reduces to the conditional MLE
$\hat\theta_n\in\arg\max_\theta\prod_i P_{Y\mid X,\theta}(Y_i\mid X_i)$.
For this reason, we develop the rest of the theory in the unsupervised setup
with a single variable $X\sim P^*$ rather than separate $X$ and $Y$, without
loss of generality.

\subsection{Reference predictors}
\label{sec:refobj}

The following quantities indexed by the model class and prior will enter the decomposition.

\paragraph{In-class KL projection.}
We define the closest in-class density to $P^*$ in the KL sense as
\begin{equation}
P^*_\Theta \;\triangleq\; P_{X\mid\theta^*(P^*)},
\qquad
\theta^*(P^*) \in \arg\!\inf_{\theta\in\Theta} D\!\left(P^*\,\big\|\,P_{X\mid\theta}\right).
\label{eq:projection}
\end{equation}
We assume the infimum in~\eqref{eq:projection} is attained.
Note $P^*=P^*_\Theta$ exactly when
$P^*$ lies in the class.

\paragraph{Bayes predictors and batch regret.}
Following \citet{polyanskiy-wu}, we define the i.i.d.\ batch redundancy
\begin{equation}
\mathrm{BatchReg}^*_{n+1}(\Theta) \;\triangleq\;
\inf_{Q_{n+1}(\,\cdot\,;\,x^n)} \,\sup_{\theta_0\in\Theta}\,
\mathbb{E}_{X^n\sim P_{X^n\mid\theta_0}}\!\Big[D\!\left(P_{X\mid\theta_0}\,\big\|\,Q_{n+1}(\,\cdot\,;\,X^n)\right)\Big].
\label{eq:batchreg}
\end{equation}
A \emph{minimax predictor} $Q^*_{n+1}$ is any minimizer of the outer infimum
(when attained). For a proper
prior $\pi$ with a well-defined posterior, the Bayes predictor is
\begin{equation}
Q^\pi_{n+1}(x_{n+1};\,x^n)
= \mathbb{E}_{\theta\sim\pi(\,\cdot\mid X^n=x^n)}
   \!\big[p_\theta(x_{n+1})\big].
\label{eq:bayesmix}
\end{equation}
Write its worst-case in-class risk and its \emph{prior cost} as
\begin{align}
\mathrm{BatchReg}^\pi_{n+1}(\Theta)
&\triangleq \sup_{\theta_0\in\Theta}
\mathbb{E}_{X^n\sim P_{X^n\mid\theta_0}}
\!\Big[D\!\left(P_{X\mid\theta_0}\,\big\|\,Q^\pi_{n+1}(\,\cdot\,;\,X^n)\right)\Big],
\label{eq:prior-risk}\\
\Gamma^\pi_{n+1}(\Theta)
&\triangleq \mathrm{BatchReg}^\pi_{n+1}(\Theta)
-\mathrm{BatchReg}^*_{n+1}(\Theta)\ge 0.
\label{eq:prior-cost}
\end{align}
Non-negativity follows because $Q^\pi_{n+1}$ is feasible in~\eqref{eq:batchreg}.
When a least-favorable prior $\pi^*$ exists whose Bayes predictor attains
the minimax value, $\Gamma^{\pi^*}_{n+1}(\Theta)=0$.

\paragraph{MLE-vs-Bayes gap.}
For $\theta_0\in\Theta$, define
\begin{equation}
\Delta^{\mathrm{MLE},\pi}_n(\theta_0)
\;\triangleq\;
\mathcal{R}_n(P_{X\mid\theta_0})
\;-\;
\mathbb{E}_{X^n\sim P_{X^n\mid\theta_0}}\!\Big[D\!\left(P_{X\mid\theta_0}\,\big\|\,Q^\pi_{n+1}(\,\cdot\,;\,X^n)\right)\Big],
\label{eq:DeltaMLE}
\end{equation}
the difference at the well-specified $\theta_0$ between the plug-in risk of MLE
and the in-class risk of $Q^\pi_{n+1}$.

\subsection{Decomposition}
\label{sec:decomp}

The plug-in risk under $P^*$ admits the following bound.

\begin{theorem}[Test-loss decomposition]
\label{thm:decomp}
For every proper prior $\pi$ satisfying the assumptions above and every $n\ge 1$,
\begin{equation}
\begin{aligned}
L_n(P^*)-H(P^*)
&\;\le\;
\underbrace{\mathcal{R}_n(P^*)-\mathcal{R}_n(P^*_\Theta)}_{\text{(I)\ approximation \& shift}}
\;+\;
\underbrace{\mathrm{BatchReg}^*_{n+1}(\Theta)}_{\text{(II)}}\\
&\quad+\;
\underbrace{\Gamma^\pi_{n+1}(\Theta)}_{\text{(II') prior cost}}
\;+\;
\underbrace{\sup_{\theta_0\in\Theta}\Delta^{\mathrm{MLE},\pi}_n(\theta_0)}_{\text{(III)\ MLE-vs-Bayes gap}}.
\end{aligned}
\label{eq:decomp}
\end{equation}
\end{theorem}
\begin{proof}
By~\eqref{eq:plugin-risk}, $L_n(P^*)-H(P^*)=\mathcal{R}_n(P^*)$. Adding and
subtracting $\mathcal{R}_n(P^*_\Theta)$ gives
\begin{equation*}
L_n(P^*)-H(P^*)
\;=\;\big[\mathcal{R}_n(P^*)-\mathcal{R}_n(P^*_\Theta)\big]
\;+\;\mathcal{R}_n(P^*_\Theta).
\end{equation*}
Since $P^*_\Theta=P_{X\mid\theta^*(P^*)}$ with $\theta^*(P^*)\in\Theta$, we have
$\mathcal{R}_n(P^*_\Theta)\le\sup_{\theta_0\in\Theta}\mathcal{R}_n(P_{X\mid\theta_0})$.
Decomposing the well-specified plug-in risk via~\eqref{eq:DeltaMLE},
\begin{equation*}
\mathcal{R}_n(P_{X\mid\theta_0})
\;=\;
\mathbb{E}_{X^n\sim P_{X^n\mid\theta_0}}\!\Big[D\!\left(P_{X\mid\theta_0}\,\big\|\,Q^\pi_{n+1}(\,\cdot\,;\,X^n)\right)\Big]
\;+\;\Delta^{\mathrm{MLE},\pi}_n(\theta_0),
\end{equation*}
and applying $\sup(a+b)\le\sup a+\sup b$ together with the definition
of $\mathrm{BatchReg}^\pi_{n+1}(\Theta)$ in~\eqref{eq:prior-risk} yields
(II) + (II') + (III).
\end{proof}

We use the right-hand side to separate candidate contributions to test loss. We now
describe each term and how it relates to double descent.

\paragraph{(I) Approximation and shift.}
The first term is the difference between the plug-in risk of MLE on the experiment
generated by $P^*$ and on the experiment generated by its in-class projection
$P^*_\Theta$. We note that this is not the textbook approximation error
$D(P^*\|P^*_\Theta)$. It need not be non-negative or monotone at finite $n$,
but vanishes when $P^*=P^*_\Theta$.

\paragraph{(II) Batch regret.}
The second term is the textbook minimax batch regret of $\Theta$. It is
monotonically non-decreasing as $\Theta$ grows: for $\Theta_1\subseteq\Theta_2$,
any $Q$ feasible in the outer infimum of~\eqref{eq:batchreg} for $\Theta_2$ is
also feasible for $\Theta_1$, while the inner supremum runs over a larger set,
so $\mathrm{BatchReg}^*_{n+1}(\Theta_1)\le\mathrm{BatchReg}^*_{n+1}(\Theta_2)$.
The additional prior cost (II') need not be monotone across model classes.

\paragraph{(III) MLE-vs-Bayes gap.}
The MLE-vs-Bayes gap does not admit a clean
asymptotic description in the non-regular settings characteristic of
deep neural networks. Our experiments find a large SGD-vs-ensemble gap
near the interpolation threshold, consistent with a large contribution
from~(III). We depict the
intuition behind this in the left panel of Figure~\ref{fig:ensemble-collapse}.

\subsection{Approximating $Q^\pi_{n+1}$ by a Gibbs ensemble}
\label{sec:gibbs}

The Bayes predictor $Q^\pi_{n+1}$ in~\eqref{eq:bayesmix} is rarely available in closed
form. However, we can approximate it by Monte~Carlo. For an inverse temperature $\beta>0$
and a prior $\pi$ on $\Theta$, the \emph{log-loss Gibbs estimator} is
\begin{equation}
\rho_{\beta,\pi}(d\theta\mid x^n)
\;\propto\; \exp\!\big\{-\beta\,\mathcal{L}_n(\theta;x^n)\big\}\,\pi(d\theta),
\qquad
\mathcal{L}_n(\theta;x^n)\triangleq \sum_{i=1}^n\log\tfrac{1}{p_\theta(x_i)},
\label{eq:gibbs}
\end{equation}
and at $\beta=1$, $\rho_{1,\pi}(\,\cdot\mid x^n)$ is exactly the Bayes posterior
under $\pi$.

For continuous parameters, we use temperature-one SGLD with unbiased
stochastic gradients of the mean negative log-likelihood and the prior drift in
Appendix~\ref{app:training}. Take
\[
\eta_t=\eta_0(1+t/\tau)^{-\alpha},\qquad
\tfrac12<\alpha\le1,\qquad
\operatorname{Cov}(\varepsilon_t)=2\eta_t I/n,
\]
where $\eta_0,\tau>0$ and $\varepsilon_t$ is the injected Gaussian noise.
Under the stability assumptions of \citet[Theorem~7]{teh2016sgld},
the step-size-weighted empirical distribution converges weakly to the posterior
as the number of steps tends to infinity.
For the sampling results below, we assume an ideal oracle that returns
independent draws from this limiting posterior.

Drawing
$\theta^{(1)},\dots,\theta^{(K)}\stackrel{\textit{i.i.d.}}{\sim}\rho_{1,\pi}(\,\cdot\mid X^n)$
and averaging the per-sample predictives gives the posterior-predictive Monte~Carlo
estimator
\begin{equation}
\widehat Q^{(K)}_{n+1}(\,\cdot\,;\,X^n)
\;\triangleq\; \frac{1}{K}\sum_{k=1}^{K} p_{\theta^{(k)}}(\cdot)
\;\xrightarrow[K\to\infty]{\text{a.s.}}\;
Q^\pi_{n+1}(\,\cdot\,;\,X^n),
\label{eq:gibbs-mc}
\end{equation}
pointwise in $X^n$ and in the test argument \citep[Ch.~6]{robert2007bayesian}.

\paragraph{Lower-bounding the MLE-vs-Bayes gap.}
Under the posterior-sampling assumption, the ensemble is pointwise
unbiased for $Q^\pi_{n+1}$ conditional on $X^n$.
Jensen's inequality applied to \eqref{eq:gibbs-mc} gives
\begin{equation}
\mathbb{E}_{\theta^{(1)},\ldots,\theta^{(K)}}\!\left[L^{\widehat Q^{(K)}}_n(P^*)\right]
\;\ge\; L^{Q^\pi}_n(P^*),
\label{eq:jensen-bound}
\end{equation}
so an unbiased predictive approximation has test log-loss at least that of
$Q^\pi_{n+1}$ in expectation over posterior sampling. Setting $\theta_0=\theta^*(P^*)$
in~\eqref{eq:DeltaMLE} under $P^*_\Theta=P^*$ gives
$\Delta^{\mathrm{MLE},\pi}_n(\theta^*(P^*))=L_n(P^*)-L^{Q^\pi}_n(P^*)$;
combining with~\eqref{eq:jensen-bound},
\begin{equation}
L_n(P^*) - \mathbb{E}_{\mathrm{MC}}\!\left[L^{\widehat Q^{(K)}}_n(P^*)\right]
\;\le\; \Delta^{\mathrm{MLE},\pi}_n(\theta^*(P^*))
\;\le\; \sup_{\theta_0\in\Theta}\Delta^{\mathrm{MLE},\pi}_n(\theta_0)
\;=\; (\mathrm{III}).
\label{eq:gap-lower-bounds-III}
\end{equation}
Thus, in the well-specified setting, the expected MLE-vs-ensemble gap
lower-bounds (III) for every ensemble size $K$ under the posterior-sampling
assumption.
